\section{The Ordinary read as one thing}

Read end to end, the Ordinary of the 1962 Missal is not the smooth object the
phrase ``the unchanging texts'' suggests. It is a deposit with visible seams,
and the seams are where its meaning is most legible.

\subsection*{What the reading has actually shown}

The exposition has moved through the rite in order, and at each stage it has
asked where the text came from as well as what it says. Setting the answers
side by side gives a result that no single movement discloses.

\begin{unittable}
\unitrow{Preface dialogue}{1043--1087}{The oldest continuously attested text in the Latin Mass, and the only one that is a negotiation rather than a prayer. Structurally unchanged.}
\unitrow{Roman Canon}{1088--1112}{Paralleled clause for clause at Milan about 390 in \work{De sacramentis} and verbally different there at every point of comparison. Trent, session XXII chapter IV, describes it as composed of the Lord's words, apostolic traditions and papal institutions.}
\unitrow{\latincue{Kyrie}, \latincue{Gloria}, \latincue{Sanctus}, \latincue{Agnus Dei}}{1024--1025, 1086, 1120}{Received chants of Greek, scriptural or early-Latin origin, added at datable moments and never composed for this rite. The \latincue{Agnus Dei} is attributed to Sergius I (687--701).}
\unitrow{Creed}{1029}{The only text in the Ordinary composed by a general council, admitted at Rome probably in 1014.}
\unitrow{Foot of the altar}{1013--1023}{Medieval private preparation, various by use until 1570, fixed thereafter, and still rubrically conditional in 1962.}
\unitrow{Offertory prayers}{1030--1042}{The youngest stratum: fourteenth century at Rome on Fortescue's reconstruction, adopted from the fourteenth Roman \work{Ordo} by the revisers of 1570.}
\unitrow{Priest's Communion prayers}{1121--1131}{Medieval private devotion, first person singular throughout, admitted with the rest in 1570.}
\unitrow{Last Gospel}{1139}{A priest's thanksgiving until 1570, when it entered the Mass; six occasions on which it is still omitted entirely.}
\unitrow{Communion of the faithful}{nn.~502--503}{Legislated in the general rubrics and not printed in the \work{Ordo Missae} at all. Its immediately preceding \latincue{Confiteor} was suppressed by the code of 1960.}
\end{unittable}

Two patterns run across that table.

The first is that \emph{the oldest material is public and the newest is
private}. The dialogue, the Canon and the chants are the Church's; the prayers
at the foot of the altar, the offertory prayers and the three prayers before
Communion are the priest's, said in the first person, and every one of them is
medieval. The rite as printed in 1962 therefore has an old public spine with a
medieval private casing at each end, and the casing is exactly the part a
congregation cannot hear.

The second is that \emph{the rubrics are younger than any of it}. Nothing in
the texts of the Ordinary was composed in the twentieth century; a great deal of
the law governing them was. The numbers this study has quoted most often ---
424, 425, 431--432, 475--476, 500--503, 507--510, 511--512 --- are the code of
1960, promulgated two years before the book that prints them. A reader who
asks what is new in the 1962 Ordinary will find the answer almost entirely in
its rubrics, and in four words of the \latincue{Communicantes}.

\subsection*{How the book makes one thing of it}

The 1962 volume does three things, none of them argumentative, that present
this stratified deposit as a single action.

It \term{numbers continuously}. The margin runs 1013 to 1139 without a break
across four printed stretches and seventy-six intervening pages of prefaces.
The Canon's own title page and running head do not restart the count. Whatever
else the numbering is, it is a claim that the reader is looking at one thing.

It \term{puts the rubric inside the prayer}. The threefold \latincue{mea culpa}
is broken by the direction to strike the breast; \latincue{Qui pridie quam
pateretur} is broken by the direction to take the host; \latincue{Supplices te
rogamus} is broken by a kiss of the altar and four signs of the cross. A rubric
set beside a prayer can be performed at the wrong moment; a rubric set inside a
sentence cannot. The effect is that no gesture in this rite is free-standing,
and no prayer is a text merely to be read.

It \term{silences almost everything}. General rubric n.~511 lists what is said
aloud at a read Mass, and between the \latincue{Sanctus} and the end of the
Canon the list contains one item. n.~512 defines the silence as audibility to
the speaker alone. The two things the people are given to hear from the
priest's own voice near the altar are \latincue{Nobis quoque peccatoribus} and
\latincue{Domine, non sum dignus}: an admission of sin and an admission of
unworthiness, both accompanied by a striking of the breast.

\subsection*{The 1962 alteration, and its size}

This edition contains one substantive change to the Canon, and it is the
cleanest available test of everything above. On 13 November 1962 the Sacred
Congregation of Rites decreed the insertion of \latincue{sed et beati Ioseph,
eiusdem Virginis Sponsi} into the \latincue{Communicantes}, effective 8
December; the facsimile collated here prints it and the 1861 hand missal does
not. Four words were added, by a pope and in the lifetime of the priests who
used the book, to the one text of which Benedict XIV had written that no pope
had added to it or changed it since Gregory the Great.

Set beside that, the change the following decade made to the same prayer's
neighbourhood is of another order. The 1962 form over the chalice reads
\latincue{Hic est enim Calix Sanguinis mei, novi et aeterni testamenti:
mysterium fidei: qui pro vobis et pro multis effundetur in remissionem
peccatorum}, and is followed by \latincue{Haec quotiescumque feceritis, in mei
memoriam facietis}. Both of those clauses were altered when the Order of Mass
was reformed. The Claude study of the postconciliar Order of Mass reaches, from
the other side of the same comparison and on its own evidence, the judgment
that the alteration of the dominical words is the deepest textual change
between the two books. Nothing in the present study contradicts that; the 1962
witness collated here is one half of the fact on which it rests, and the size
of the 1962 insertion is the measure against which the later change should be
read.

\subsection*{Project synthesis}

\begin{studybox}{Project synthesis --- the judgment this study reaches}
\textbf{The judgment.} The Ordinary of the 1962 \work{Missale Romanum} is a
stratified text whose strata are separately datable and demonstrably unlike
each other in age, authorship and register; its unity is not the unity of a
single composition and never was. What the 1962 book supplies is not that
unity but its \emph{exhibition}: a continuous numbering, a rubric embedded
inside every prayer, and a rule of silence that reduces the audible Mass to a
few lines, together present the deposit as one ordered action directed to one
end. The rite's identity therefore rests on what it does and on the law that
orders it, not on the homogeneity of its texts --- and the 1962 books are
candid about this, because their own promulgating council said the Canon was
made of three different kinds of material and their own Congregation of Rites
added four words to it while the edition was in the press.

\smallskip
\textbf{The reasons.}
\begin{itemize}[itemsep=0.15em]
\item \emph{The strata are visible in the sources, not inferred from a theory.}
\work{De sacramentis} IV, 5--6, read at the page images, has
\latincue{ascriptam, ratam, rationabilem, acceptabilem} and \latincue{memores}
with the same three genitives, and differs from the Roman text at every point
of comparison; Fortescue dates the offertory prayers to the fourteenth century
at Rome and the last Gospel's admission to 1570; Berno's report dates the Creed
at Rome to 1014; the \work{Liber Pontificalis} attributes the \latincue{Agnus
Dei} to Sergius I. These are four independent kinds of evidence pointing the
same way.
\item \emph{The register changes with the stratum.} The old public material is
plural and addressed to the Father; the medieval material is singular and
addressed to Christ or to the Trinity. \latincue{Perceptio Corporis tui \dots\
non mihi proveniat in iudicium} and \latincue{Te igitur, clementissime Pater
\dots\ supplices rogamus ac petimus} are not two moods of one voice.
\item \emph{The unifying devices are the book's own and are checkable.} The
numbering runs 1013--1139 unbroken; nn.~511--512 define the silence and its
two exceptions; the interleaved rubric appears in every appointed block quoted
in this study. None of these is an editorial construction.
\item \emph{The law is doing the work that a single authorship would otherwise
have to do.} The rubrics decide when the psalm is said, when the
\latincue{Gloria} and Creed are said, which preface is used, whether there is a
last Gospel, and whether the people's \latincue{Confiteor} is said before their
Communion. Change the code and the celebration changes without a word of the
Ordinary changing --- which is exactly what happened in 1960.
\item \emph{The 1962 insertion is a controlled case.} Four words entered the
Canon by decree, and the two witnesses collated for this study disagree about
them and about nothing else of substance. That the rite was not thought to have
changed is evidence about where its identity is located.
\end{itemize}

\smallskip
\textbf{The strongest counterargument, at full strength.} The account above is
an archaeologist's, and a rite is not an archaeological object. Every serious
liturgical text is stratified --- so is the Divine Office, so is every Eastern
anaphora, so is the Bible --- and pointing this out establishes nothing about
identity. The Church itself says the Canon is made of the Lord's words,
apostolic traditions and papal institutions, and says so in the act of
declaring it free from error; on the Catholic account the composite origin is
not a tension to be explained but the ordinary manner in which a living rite
grows. Worse, by locating the unity in the 1962 book's numbering, typography
and rules of voice, this study has made an accident of one printing into the
principle of the rite. Priests who used the typical editions of 1884 or 1920
had different pagination, no continuous marginal numbering of this kind and a
different code of rubrics, and they celebrated the same Ordinary; a Dominican
or a Carmelite of the same date, with a different \latincue{Confiteor}, a
different offertory and different rubrics, was also offering the Roman Mass in
the sense that matters. If the identity of the rite survives all that, then it
cannot be constituted by the devices this study has named, and the real unity
of the Ordinary is what the counterargument says it is: a single sacrificial
act, whose texts are its instruments and whose printing is merely how one
century handed it on. On that view the correct description is the simple one
--- the Ordinary is the Mass, stated in words --- and the stratigraphy is a
footnote.

\smallskip
\textbf{Why the judgment is nevertheless retained.} Because the counterargument
and the judgment are answering different questions, and the counterargument's
answer is not available to a study that has bound itself to one identified
book. That the Roman Ordinary is one sacrificial act is a theological claim
which this study neither doubts nor can establish from a facsimile; what a
facsimile can establish is how this edition presents that act, and the answer
is: by these three devices and by this code. The judgment above is deliberately
about exhibition rather than constitution, and it says so. It also survives the
counterargument's best example in the one form in which this study can test it.
The 1861 hand missal, printed from a typical edition three-quarters of a
century older than this one, already sets its rubrics inside the prayers at the
same places --- inside the fraction, inside the last Gospel at \latincue{Et
Verbum caro factum est} --- and already prints the same conditional omissions
at Requiems. The devices are therefore not an accident of this printing, even
though their 1962 form is peculiar to it. And the medieval variants Fortescue
records --- Salisbury, Dominican, Carthusian, Carmelite, Lyons, all of them at
the prayers at the foot of the altar --- fall exactly in the stratum this study
has identified as the youngest and most variable. The counterargument, pressed
on its own examples, turns out to support the stratification it dismisses.

\smallskip
\textbf{What specific finding would change it.} One would be decisive:
manuscript evidence that the private medieval strata are not later additions to
a public Roman core but were present in Roman books contemporaneously with the
Canon --- concretely, a Roman \work{Ordo} or sacramentary earlier than the
eleventh century containing \latincue{Suscipe, sancte Pater} or the sequence
\latincue{Aufer a nobis}--\latincue{Oramus te} in the position the 1962 book
gives them. Fortescue's reconstruction, on which this study's dating of the
offertory depends, rests on manuscript arguments this study has not examined,
and it is the load-bearing external claim in the whole account. If the
offertory prayers proved coeval with the Canon, the ``old public spine, young
private casing'' description would collapse, the register argument would lose
its chronological force, and the judgment would have to be rewritten as a claim
about literary genre rather than about strata. A second, weaker finding would
also bite: evidence that the continuous marginal numbering of 1013--1139 was
introduced for typesetting convenience with no editorial intention behind its
continuity, which would remove one of the three devices named above without
touching the other two.

\smallskip
\textbf{Status.} This is a source-grounded editorial synthesis by an AI
contributor, resting on one collated facsimile of the 1962 typical edition, one
collated public-domain English hand missal of 1861, and the identified
historical and conciliar witnesses named in the references. It is not
magisterial teaching, a canonical opinion, a critical edition, or a specialist
consensus, and it has received no independent liturgical, historical or
theological review.
\end{studybox}

\subsection*{What this study does not decide}

It does not date the sections of the Roman Canon severally, assign an author to
any of them, or adjudicate the argument about the Roman Canon's want of an
explicit epiclesis. It does not identify the \latincue{sanctus Angelus tuus} of
\latincue{Supplices}. It does not decide whether the medieval offertory prayers
improved or damaged the rite. It offers no account of the present canonical
discipline governing the use of the 1962 books, which is a dated question of
current law and not a question about this edition. And it makes no comparison
between this Ordinary and the reformed Order of Mass beyond the two textual
points at which the collation of this edition supplies evidence: the four words
added to the \latincue{Communicantes} in 1962, and the wording of the two
dominical forms as this book prints them.
